Una matriu quadrada $A$ és \emph{idempotent} si $A^2=A$.

\begin{apartats}

\apartat[1,25]{0,75}
Comprova que $A=\begin{pmatrix}2&-2\\1&-1\end{pmatrix}$ és idempotent, i dedueix-ne $A^n$ per a
qualsevol nombre natural $n$.

\begin{solucio}
\[
A^2=\begin{pmatrix}4-2&-4+2\\2-1&-2+1\end{pmatrix}=\begin{pmatrix}2&-2\\1&-1\end{pmatrix}=A.
\]
Aleshores $A^3=A^2A=AA=A^2=A$ i, repetint-ho, $A^n=A$ per a qualsevol $n\ge1$.
\end{solucio}

\begin{tria}{tipus-potencies}
\itemtria{parametre-idempotent}{1}{1,25}
Troba el valor de $m$ perquè la matriu $M=\begin{pmatrix}m&2\\-1&-1\end{pmatrix}$ sigui idempotent.

\begin{solucio}
\[
M^2=\begin{pmatrix}m^2-2&2m-2\\-m+1&-1\end{pmatrix}.
\]
Cal que $M^2=M$, element a element: $m^2-2=m$, $2m-2=2$, $-m+1=-1$ i $-1=-1$. La primera equació dona
$m=2$ o $m=-1$; la segona i la tercera, $m=2$. Només $m=2$ les compleix totes: $M$ és idempotent si i
només si $m=2$.
\end{solucio}

\itemtria{periodica}{1}{1,25}
Sigui $B=\begin{pmatrix}0&-1\\1&0\end{pmatrix}$. Calcula $B^2$, $B^3$ i $B^4$, i fes-los servir per
trobar $B^{2027}$.

\begin{solucio}
\[
B^2=\begin{pmatrix}-1&0\\0&-1\end{pmatrix}=-I,\qquad
B^3=B^2B=-B=\begin{pmatrix}0&1\\-1&0\end{pmatrix},\qquad
B^4=(B^2)^2=I.
\]
Les potències es repeteixen cada 4. Com que $2027=4\cdot506+3$,
$B^{2027}=(B^4)^{506}B^3=B^3=\begin{pmatrix}0&1\\-1&0\end{pmatrix}$.
\end{solucio}
\end{tria}

\begin{nomesllarg}
\apartat{0,75}
Demostra que, si una matriu $A$ és idempotent, la matriu $I-A$ també ho és.

\begin{solucio}
$(I-A)^2=I-A-A+A^2=I-2A+A$, perquè $A^2=A$. Queda $I-A$: és idempotent. (La identitat commuta amb $A$,
i per això el quadrat es desenvolupa com amb nombres.)
\end{solucio}
\end{nomesllarg}

\end{apartats}
